\section{電気探査：地下へ流した電流から何が分かるか}
\label{theory:exploration:dcarrays}
地層を掘る前に、水や粘土、空隙、金属などによる電流の流れやすさの差を調べたい。比抵抗法は電流を注入して地表や孔内の電位差を測り、地下の抵抗率分布を推定する。抵抗\(R\)は試料の形状も含むが、均質な棒では\(R=\rho L/A\)である。長さ\(L\)、断面積\(A\)から\(\rho=RA/L\)となるため、抵抗率の単位は\(\Omega\,\mathrm m\)であり、\(\Omega\,\mathrm m^3\)ではない。局所形では\(\boldsymbol E=\rho\boldsymbol J\)で、導電率は\(\sigma=1/\rho\)である。

定常電流の保存と\(\boldsymbol E=-\nabla V\)を合わせると、源のない部分では
\[
 \nabla\cdot(\sigma\nabla V)=0
\]
となる。注入電極を源・回収電極を吸込みとして、地表の境界条件と地形を含めて電位を計算するのが順問題である。多数の電位差へ一致する\(\sigma\)を探すのが逆問題で、電位だけで地質名が一意に決まるわけではない\cite{epa_resistivity}。

\subsection*{半空間の電位とWennerの幾何係数}
教材の説明用に均質・等方な導電半空間、絶縁した平らな地表、小さい点電極を仮定する。地表の注入点から半球面を通る電流を\(I\)とすると、\(J_r=I/(2\pi r^2)\)である。したがって無限遠基準の電位は\(V(r)=\rho I/(2\pi r)\)となる。注入Aと回収Bの重ね合わせより、電位電極M,Nの差は
\[
 \Delta V=\frac{\rho I}{2\pi}
 \left(\frac1{AM}-\frac1{BM}-\frac1{AN}+\frac1{BN}\right).
\]
Wenner配置を地表直線上のA--M--N--B、全隣接間隔\(a\)とすると、括弧は\(1/a\)であり
\[
 \rho_a=2\pi a\frac{\Delta V}{I}
\]
を得る。この導出は教材独自である。符号は電極と電位差の定義をそろえる。不均質な地盤の\(\rho_a\)は「この配列で同じV/Iを出す均質半空間」の値であり、MとNの真下にある一層の真の抵抗率ではない。

Schlumberger、dipole--dipole等も幾何係数と地下感度が異なる。垂直探査では中心をおおむね固定し間隔を広げ、水平多層等のモデルを推定する。水平探査では配置を移動して横の変化を調べる。固定間隔でも感度は有限体積に分布するので「特定の深度だけの水平抵抗率」を直接測ると読まない。二次元・三次元の比抵抗トモグラフィは、多数配置のデータを順モデルへ合わせて分布を求める。偽断面の表示位置を実深度と同一視しない。

\subsection*{線形フィルタは配置でなく順応答の計算法}
層状大地モデルの比抵抗変換から見掛比抵抗曲線を計算する際、線形フィルタを用いる方法がある。USGSのTI-59プログラムはその具体例である\cite{usgs_linear_filter1981}。平成23年度問30では、配列名に混じるリニアフィルタを一般の雑音除去だけと解釈しない。IPは追加の分極応答を測るが、比抵抗と同じ電極を使うため、無関係という全称も避ける。

\subsection*{四電極で分ける接触電圧降下}
\label{theory:exploration:fourpoint}
二端子で電流を流しその端子電圧を測ると、地盤だけでなく両電極の接触抵抗による電圧降下も含む。別のM,Nを高入力抵抗\(R_{\rm in}\)の電圧計へつなぎ、測定電流\(I_v\)を十分小さくすれば、接触誤差の代表値\(I_v(R_{cM}+R_{cN})\)を小さくできる。これが四電極の要点である。送信側の接地抵抗は電流を確保する問題として残り、四電極なら接地・地形・電磁カップリングの影響が全てゼロになるわけではない。平成25年度問34はこの分離を問う。

\section{含水岩石の電流経路とArchie式}
\label{theory:exploration:archie}
多くの岩石骨格は絶縁性が高く、孔隙水に溶けたイオンが電流を運ぶ。孔隙率\(\phi\)、水飽和度\(S_w\)、水の抵抗率\(\rho_w\)について、代表的な実験モデルは
\[
 \rho_t=a\rho_w\phi^{-m}S_w^{-n}
\]
である。\(a,m,n\)は骨格・連結性・濡れ条件等に依存する係数で、全岩石に同じ値を置かない。飽和では\(\rho_t/\rho_w=a\phi^{-m}\)が地層係数となる。孔隙ネットワークの少ない導電断面と迂回路をまとめたモデルであり、普遍的な保存則から全係数が決まる式ではない。

粘土表面、導電鉱物、電子伝導に無視できない電流があれば、孔隙水のみという前提が破れる。単純な並列伝導モデル\(\sigma_{\rm bulk}=\sigma_{\rm pore}+\sigma_{\rm surface}\)を考えると、追加の表面伝導は抵抗率を下げる。低抵抗の水は相対的に孔隙水伝導を優勢にしやすい。逆に高抵抗水だけでは式の精度を保証できない。ただし清浄な絶縁骨格なら淡水でも利用し得る\cite{epa_borehole_resistivity}。

海水の塩分は通常導電率を上げ抵抗率を下げる。不飽和化は連続水経路を減らし高抵抗化しやすい。粘土化・含水割れ目は低抵抗化しやすいが、乾燥した割れ目や珪化まで同じ傾向にはしない。「礫層は高抵抗」という典型も周辺の粘土層との対比・塩分条件による。原典の一般傾向を、岩石名から値を一意に決める規則にしない。

\section{IPとSP：人工分極と自然電位の機構}
\label{theory:exploration:ipsp}
IP（誘導分極）は人工電流を流した際の地下界面の電荷蓄積を利用する。電流遮断後の電位減衰又は周波数応答の位相を測る。硫化鉱物の電子伝導と電解質のイオン伝導の界面、鉱物表面等が応答へ寄与する。比抵抗だけでは区別しにくい界面状態を調べるための追加情報であり、IPでも比抵抗を測る。ケーブル間の電磁結合が分極応答へ混ざり得るため、周波数・配置・測定時間を吟味する\cite{epa_ip}。

SP（自然電位）は人工的に電流を注入せず、基準点との電位差を非分極電極で測る。地下水流による電気二重層の電荷輸送、酸化還元差を結ぶ導電性鉱体、温度や化学勾配等が源となる。説明用に\(\boldsymbol J=-\sigma\nabla V+\boldsymbol J_s\)と置けば、電流保存は
\[
 \nabla\cdot(\sigma\nabla V)=\nabla\cdot\boldsymbol J_s
\]
となる。観測から源\(\boldsymbol J_s\)を推定するには抵抗率と源機構のモデルが必要である。電位の絶対零点は測らず、基準点の選び方も記録する。

典型的な硫化鉱体の鉱化電池では地表の負異常を利用する。一方、流動・熱・化学が重なる場の符号は一つの法則へ固定できない。正異常だから必ず熱水、負異常だから必ず鉱体と断定しない。SP異常の幅や幾何モデル、他観測との同時逆解析による深度推定は可能だが、SP分布だけで深度が一意に読めるわけではない\cite{epa_sp}。令和3年度問32の公式4は保持するが、「SPでは深度解析ができない」という誤った禁止を教材にしない。原典の否定範囲に残る疑義を問題ごとに示す。

\section{有限送信源・周波数領域・時間領域}
\label{theory:exploration:nearfield}
自然MTと人工CSAMTの抵抗率・表皮深度は前節\ref{theory:exploration:electromagnetic}の式へつながる。均質導電半空間で変位電流を無視すると、周期場の拡散による振幅減衰尺度は\(\delta=\sqrt{2\rho/(\omega\mu)}\)である。送受信距離\(r\)に対して\(r/\delta\)が十分大きい遠方では平面波近似を使いやすい。固定rで低周波化又は高抵抗化すればδが増し、近接場へ移る。USGS実地報告は三表皮深度超を目安として用いるが、その数を全装置・全構造の厳密閾値とはしない\cite{usgs_csamt2012}。

\(\rho_a=|E/H|^2/(\omega\mu)\)のµは透磁率である。真空ではµ0を用い、誘電率ε0ではない。HをBと取り違えると換算係数と単位も変わる。平成28年度問23は式の形だけでなく、この変数定義の誤りを問う。

\subsection*{周波数領域では振幅と位相を測る}
\label{theory:exploration:fdtd}
FDEMでは正弦電流が送信コイルの一次磁場を作り、地盤の渦電流が二次磁場を生む。定常応答の複素振幅比を同相・直交成分又は振幅・位相として測る。Slingramやloop--loopはこの仕組みを使う。電極を挿さなくてもよいが、コイル間隔、周波数、導電率、姿勢に感度が依存する\cite{epa_fdem}。

TDEM/TEMでは送信を遮断した後の過渡応答を測る。均質導電体の準静的近似では\(\partial\boldsymbol B/\partial t=(1/(\mu\sigma))\nabla^2\boldsymbol B\)となるため、代表拡散距離は\(\sqrt{t/(\mu\sigma)}\)の尺度で増す。遅い時刻には一般に深部への感度があるが、信号は弱くなり、装置寸法・電流・遮断時間・測定時間・ノイズが限界を決める。central loopで深部を測れることを「送受信間隔によらず無制限に深く測れる」としない\cite{epa_tdem_archive}。一次場の遮断後観測と、電場の局所歪みによるMTのstatic shiftの少なさは別の性質である。

VLFは既存通信局の3--30kHz帯の信号を使う代表的方式で、局の方位と対象構造の向きが重要である。調査者は既存局の送信条件を任意には変更できない。Schumann共振も電気電磁現象であり、MT観測記録で調べられている\cite{schumann_mt2009}。反射地震のNMO補正とは分野が異なる。

\section{GPR：走時断面から電気境界を推定する}
\label{theory:exploration:gpr}
地中レーダは高周波電磁パルスを放射し、電気特性の異なる境界からの反射を受ける。低損失・非磁性の均質媒質なら
\[
 v\simeq\frac{1}{\sqrt{\mu_0\epsilon}}
 =\frac{c}{\sqrt{\epsilon_r}},\qquad
 \epsilon_r=\frac{\epsilon}{\epsilon_0},\qquad
 z\simeq\frac{vt}{2}
\]
である。tは往復走時であり、直下水平面などの単純な幾何と速度が必要である。εはF/m、εrは無次元で、空気を分母にするのは通常の近似に過ぎない。高導電性・分散媒質ではこの単純速度式の条件が破れる。

波形をアンテナ位置へ並べるprofileはまず時間断面である。wide-angleや共通中点配置から速度を推定し、深度変換に用いる。埋設管等を横切る測線は位置を捉えやすい。点対象の回折双曲線をそのまま地層形状と思わず、速度・偏位・必要なmigrationを考える\cite{epa_gpr}。

\subsection*{分解能力と最小識別距離を区別する}
\label{theory:exploration:gprresolution}
二つの反射を時間差\(\Delta t\)で識別するにはパルス幅・帯域が制約になる。説明用には帯域Bで\(\Delta t\sim1/B\)、距離分解尺度\(\Delta z\sim v/(2B)\)である。帯域が中心周波数に比例する同系統装置なら、高周波ほど最小識別距離が小さく、能力としての「分解能」は良い。この語の二つの使い方で比例の向きが変わるので、式と変数を明示する。

実際の探査深度は導電減衰、散乱、幾何拡がり、アンテナ帯域、検出閾値等で決まる。高周波で分解を良くすると到達深度が制限されるという交換条件はあるが、深度が厳密に1/fへ比例する普遍式ではない。低損失非分散媒質の導電損失近似は\(\alpha\simeq(\sigma/2)\sqrt{\mu/\epsilon}\)で、fだけの反比例を導くものではない。

通常のGPRはMHz以上のパルスで、原典の10MHz程度は分類目安である。10kHzとの三桁差は誤りだが、10MHzに絶対的な物理境界があるわけではない。力学強度を直接観測しない一方、含水・空隙・劣化等の共通要因による間接相関はあり得る。「ほとんど相関なし」という原典はその限定を保持する。数mの空洞や埋設物の用途を、全媒質の到達上限へ変えない。

\section{孔内検層：深度目盛と観測物性を結び付ける}
\label{theory:exploration:borehole}
地表探査だけでは深度の曖昧さが残る。ボーリング孔に沿う連続計測は、その位置の深度と物性応答を結び付ける。ただし孔周辺の有限体積を見ており、孔壁・泥水侵入・ケーシング・装置結合の補正がある。検層記録の深度が分かることと、地盤全体の構造が完全に分かることは別である。

速度検層のdownholeは地表起振・孔内受振、upholeは孔内起振・地表受振である。走時から速度を求めるには距離と伝播経路、屈折・傾斜の条件を扱う。孔壁へ密着させるdownhole受振器なら水位のない部分でも測定できる\cite{geebec_ps}。水位以深だけという限定は誤りだが、孔内状態によらず精度が保証される意味ではない。

suspensionは孔内振源と複数受振器を持つゾンデで、受振点間の時間差から区間速度を得る\cite{oyo_suspension}。地盤のP/S応答を測り、水中をS波が伝わるという説明にはしない。VSPでは通常の地表振源・孔内受振と逆配置等を区別し、反射記録と孔内深度・地層を対比する。初動速度だけの検層と反射の解析目的は同じではない。

normal電気検層は電極の注入・電位差による比抵抗計測である。導電水・泥水を通じて地層へ電流を流す通常の裸孔条件を考える。孔内水の有無・抵抗率を無視しない。一方、コイルの誘導検層には別の接触条件があり、電極型の水位条件を全検層へ拡張しない。P波速度は弾性率と密度に関わり、含水・飽和で変わり得るが、含水比だけで全地盤の増減が決まるわけではない。

\subsection*{ガンマ密度：粒子密度でなくかさ密度}
\label{theory:exploration:gammadensity}
通常の密度検層はガンマ線を放射し、主にCompton散乱を介する応答から電子密度を測る。一般的岩石の組成では校正によりかさ密度へ結び付く。ベータ線やRayleigh散乱へ置き換えない。飽和二成分モデルなら
\[
 \rho_b=(1-\phi)\rho_m+\phi\rho_f,\qquad
 \phi=\frac{\rho_m-\rho_b}{\rho_m-\rho_f}
\]
である。\(\rho_m\)は粒子・骨格、\(\rho_f\)は孔隙流体の密度であり、これらが既知という条件から孔隙率を推定する。粒子密度そのものの直接測定ではない。孔壁接触や孔内流体、線源・検出器校正が必要である\cite{epa_gamma_density}。

\section{時間差・幾何・熱を使う補足調査}
\label{theory:exploration:multimethod}
time lapseは同じ地点・方式の反復測定を比較し、地下水移動や物性変化を推定する。差分には機器ドリフトや測線ずれ・温度差も入り得る。tomographyの覆いが不足すれば、特定方向の構造や薄い領域を分解しにくい。重力では同じ形状深度なら密度差の大きさが異常を増すが、密度差だけで異なる対象の検出難易度を順位付けしない。

1m深地温探査は、地表下1mの温度を測り流動水・熱源による平常温度との差を調べる。1mは測温深度であり、水脈の最大深度ではない。水温との温度差が小さい季節には検出しにくく、地況・境界条件を扱う熱モデルと観測を組み合わせる\cite{naro_thermal2008,landslide_thermal2007}。特定地の適期を全国共通の月にしない。

環境振動の振動レベルは人の感覚に応じた周波数補正を施す加速度レベルである。原典のJIS準拠測定を、単なる最大変位測定へ取り違えない\cite{rion_vibration2016}。出題時と現行の規格本文の版差は未照合であり、現在版の番号や性能要件を推測して補わない。

\subsection*{試験での確認順序}
励起源、受信器、直接観測量、推定物性、単位、有限距離・帯域・媒質の条件を順に照合する。今回の29問での論点を全16試験回の頻度の確定値とはしない。SPの深度推定、Archieの表面伝導、GPRの比例表現等の疑義は各問題へ残し、公式正答と科学的な説明の保証範囲を区別する。
